\subsection{An explicit projector for evolution factors}
For an evolution operator with $P=\prod_{j=1}^N(\tau-a_j(\xi))$ and distinct $a_j$, a component is an eigenfunction of $\partial_t$: $\partial_tu_j=a_j(\partial_x)u_j$. Applying $\prod_{k\ne j}(\partial_t-a_k(\partial_x))$ to $u=\sum_lu_l$ kills every $u_l$ with $l\ne j$ (it contains the factor $\partial_t-a_l$) and multiplies $u_j$ by $\prod_{k\ne j}(a_j-a_k)(\partial_x)$. Hence
\begin{equation}\label{xeq:lagrange}
u_j=\prod_{k\ne j}\frac{\partial_t-a_k(\partial_x)}{a_j(\partial_x)-a_k(\partial_x)}\;u ,
\end{equation}
the Lagrange formula of Eq.~\eqref{eq:spectral} with $C$ replaced by $\partial_t$. The result is a Fourier multiplier in $x$ that is rational in $\xi$, with poles where $a_j(\xi)=a_k(\xi)$; this matches the characterization ``projectors rational in $\xi$ iff the factors are coprime'' \led{II-T2/T3}.

\subsection{Where decomposition fails}
\begin{proposition}[{(Holes) \led{C12-P10} \proved[]}]
On a domain with a hole, the decomposition fails for some $P$.
\end{proposition}
Figure~\ref{fig:hole} shows the standard obstruction. In the theorem the hypothesis ``simply connected'' means $H^1_{dR}(\Omega)=0$, exactly what removes this obstruction; a cohomological criterion for non-simply connected $\Omega$ is also available \led{OP-FUN, C12-Q6}.
\begin{figure}[!htb]
\centering
\begin{tikzpicture}
\fill[cA!10] (0,0) circle (1.7); \fill[white] (0,0) circle (0.5);
\draw[cA,thick] (0,0) circle (1.7); \draw[cA,thick] (0,0) circle (0.5);
\node at (0,0) {\small hole};
\draw[cB,thick,-{Stealth}] (1.1,0) arc (0:340:1.1);
\node[cB,font=\scriptsize] at (0,-2.05) {$\oint d\arg z=2\pi$};
\node[align=left,font=\small,anchor=west] at (2.4,0) {$u=\log|z-\alpha|^2$ solves $\Delta u=0$ on the annulus,\\ but $u\neq h(z)+\overline{g(z)}$ with $h,g$\\ single-valued holomorphic: the logarithm\\ has monodromy (winding number $\neq0$).};
\end{tikzpicture}
\caption{The Laplacian $\partial_x^2+\partial_y^2=(\partial_x+i\partial_y)(\partial_x-i\partial_y)$ factors, but the factor decomposition fails on a punctured disc: the obstruction is topological ($H^1_{dR}\neq0$).}\label{fig:hole}
\end{figure}

\subsection{Regularity: when are the components in $H^m$?}
The decomposition is an algebraic statement; analysis enters when we want components in a Sobolev space. For two components the governing condition is geometric: the two component spaces $M_1$ and $M_2$ must be separated by a positive angle (Figure~\ref{xfig:friedrichs}); the next theorem makes this precise for any number of components. Two difficulties appear along the way; they are collected in the box after the theorem.
\begin{proposition}[{(Loss of a derivative) \led{C12-C1} \falsest}]\label{prop:loss}
``If $u\in H^m(\Omega)$ then all $u_j\in H^m$'' is false: one derivative can be lost. In other words, the pieces of a smooth function may be rougher than the function itself.
\end{proposition}
\begin{theorem}[{(Regularity criterion) \led{C12-T3, C12-P7} \proved[]}]\label{xthm:hm}
All components lie in $H^m$ for every $u\in\ker_{H^m}P$ iff the sum $\sum_j\ker_{H^m}Q_j$ is \emph{closed} and \emph{dense} in $\ker_{H^m}P$. Closedness is equivalent to $\gamma_N>0$, where
\[
\begin{aligned}
\gamma_N&=\inf\sigma\big(S^*S\big|_{(\ker S)^\perp}\big),\\
S&:\bigoplus_{j=1}^N\ker Q_j\to H^m,\qquad S(u_1,\dots,u_N)=\sum_{j=1}^Nu_j .
\end{aligned}
\]
For $N=2$ this is the condition $c_F<1$ on the Friedrichs cosine. Density holds on every bounded convex $\Omega$, for every $P$ and $m$ \led{OPF5B-T1} \modulo{H\"ormander~\cite{ho}, Stein}.
\end{theorem}
\begin{notebox}{Reading $\gamma_N$}
\begin{itemize}[nosep]
\item The space $\bigoplus_j\ker Q_j$ consists of the $N$-tuples $v=(u_1,\dots,u_N)$ with norm $\|v\|^2=\sum_j\|u_j\|_{H^m}^2$;
\item $Sv=\sum_ju_j$ is their sum;
\item $S^*$ is the adjoint of $S$;
\item $(\ker S)^\perp$ removes the tuples that sum to zero (the ambiguity of Theorem~\ref{thm:factor});
\item $\inf\sigma$ is the bottom of the spectrum.
\end{itemize}
Since $S^*S$ restricted to $(\ker S)^\perp$ is self-adjoint and nonnegative, the bottom of its spectrum is the minimum of the Rayleigh quotient
\[
\frac{\langle S^*Sv,v\rangle}{\|v\|^2}=\frac{\|Sv\|^2}{\|v\|^2}.
\]
Hence the formula of the theorem is equivalent to the more intuitive
\[
\boxed{\ \gamma_N=\inf_{\substack{v\in(\ker S)^\perp\\ v\ne0}}\frac{\|Sv\|_{H^m}^2}{\|v\|^2}\ }
\]
\end{notebox}

\begin{aside}{In words}
$\gamma_N$ is the worst case, over all tuples of pieces, of the ratio between the size of the sum and the size of the pieces.
\begin{itemize}[nosep]
\item If $\gamma_N>0$ the pieces cannot nearly cancel each other, and
\[
\|v\|\le\gamma_N^{-1/2}\|Sv\|:
\]
the pieces are controlled by the sum, so recovering them from $u$ is stable.
\item If $\gamma_N\to0$ there are large pieces whose sum is tiny, and the recovery is unstable.
\end{itemize}
\textbf{The case $N=2$.} One finds $\gamma_2=1-c_F$, where
\[
c_F=\sup\{|\langle a,b\rangle|:\ a\in M_1\ominus(M_1\cap M_2),\ b\in M_2\ominus(M_1\cap M_2),\ \|a\|=\|b\|=1\}
\]
is the Friedrichs cosine: the part of $v$ lying in $M_1\cap M_2$ contributes the ratio $2$ and does not affect the infimum, and on the rest
\[
\|a+b\|^2\ge(1-c_F)(\|a\|^2+\|b\|^2)
\]
with near-equality attained. Hence $c_F<1\iff\gamma_2>0$, as stated.
\end{aside}
\begin{figure}[!htb]
\centering
\begin{tikzpicture}[scale=1.1]
\draw[->] (-0.3,0)--(3.6,0); \draw[->] (0,-0.3)--(0,2.8);
\draw[cA,very thick,->] (0,0)--(3.2,0.5) node[right,font=\small]{$M_1=\ker Q_1$};
\draw[cB,very thick,->] (0,0)--(1.9,2.4) node[above,font=\small]{$M_2=\ker Q_2$};
\draw[cG] (0.9,0.14) arc (8.9:51.6:0.9); \node[font=\small] at (1.12,0.55) {$\theta$};
\node[font=\small,anchor=west] at (4.3,2.0) {$c_F=\cos\theta$};



\end{tikzpicture}
\caption{Friedrichs angle $\theta$ between two component spaces $M_1=\ker Q_1$ and $M_2=\ker Q_2$. A larger angle means a better-conditioned recovery of the pieces from their sum.}\label{xfig:friedrichs}
\end{figure}
For $N=2$, Figure~\ref{xfig:friedrichs} reads as follows: $c_F=\cos\theta$; if $\theta>0$ the sum $M_1+M_2$ is closed and the pieces depend continuously on $u$; as $\theta\to0$ the split becomes unstable. For $N\ge3$ the pairwise angles do not suffice; $\gamma_N$, the smallest nonzero eigenvalue of $S^*S$, is the right conditioning number: it is the conditioning of the linear map ``components $\mapsto$ sum'', exactly like the smallest nonzero singular value in a least-squares problem.

\begin{caution}[title=Two problems with regularity]
\textbf{Problem 1: stability \falsest.} Algebraically the split always exists (Theorem~\ref{thm:factor}), but the pieces may be less regular than $u$. When the angle between the component spaces tends to $0$ (that is, $\gamma_N\to0$) the map ``$u\mapsto$ pieces'' becomes unstable, and one derivative can be lost (Proposition~\ref{prop:loss}). For the algorithm: check $\gamma_N>0$ before using the components in $H^m$ (Step~5).\\[3pt]
\textbf{Problem 2: three or more components \openq.} For $N\ge3$ the pairwise cosines $c_F(i,j)$ do not decide whether the sum of kernels is closed: an abstract Hilbert-space counterexample exists \led{S-CHECK-4}, and whether it can occur for differential operators is open \led{C12-Q2}.
\end{caution}

\begin{takeaway}
Algebraically the split always exists (Theorem~\ref{thm:factor}); analytically, the components are well behaved exactly when the component spaces are separated by a positive angle (Theorem~\ref{xthm:hm}). The next section stops asking for a split into factors and asks what survives when one splits by additive terms instead.
\end{takeaway}

%======================================================================
